%\newcommand{\imprimirresumo}{
\begin{resumo}
\hspace{1.5cm}
Trata-se de um apanhado conciso do conteúdo do trabalho, não uma simples enumeração de tópicos. Deve conter até 500 palavras, em um único parágrafo, seguido das palavras-chave representativas do conteúdo do trabalho, conforme formato exemplificado abaixo. A elaboração do resumo é regida pela norma ABNT NBR 6028:2003a (Resumos).
\vspace{\onelineskip}
\noindent\\
\textbf{Palavras-chave}: Palavra 1. Palavra 2. Palavra 3. Palavra 4. Palavra 5.

%\textcolor{red}{(Use até 5 palavras-chave)}
\end{resumo}

\newpage

\begin{resumo}[Abstract]
\begin{otherlanguage*}{english}
\hspace{1.5cm}
Consiste na tradução do resumo para uma língua estrangeira, mantendo-se a mesma formatação e conteúdo. As línguas mais comuns são o inglês, francês e espanhol. Recomenda-se utilizar um serviço especializado de tradução, evitando-se o uso de um tradutor automático.
\vspace{\onelineskip}
\noindent\\
\textbf{Keywords}: Word 1. Word 2. Word 3. Word 4. Word' 5.
\end{otherlanguage*}
\end{resumo}

%}


